% SHARD-ID: CA-63-QUBIT-MOMENT-STATE-EXISTENCE
% SHARD-TITLE: All-Level Moment Feasibility Gives a Qubit Reference State
% SHARD-SUMMARY: Proves that exact feasibility of the monotone full-window qubit moment hierarchy yields a genuine state on the quasi-local spin algebra.
% SHARD-SUMMARY: Shows that the all-level relation constraints vanish as bounded operators in the GNS representation.
% SHARD-SUMMARY: Separates this compactness result from uniqueness, ground-state, clustering, continuum-limit, and Virasoro claims.
% SHARD-KEYWORDS: moment hierarchy, compactness, state space, GNS, qubit chain, residual quotient

\section{All-Level Moment Feasibility Gives a Qubit Reference State}
\label{sec:qubit-moment-state-existence}

\subsection{Status, Sources, and Dependencies}

\textbf{Status: derived compactness theorem, conditional on standard
C*-algebra facts.}  This upgrades CA-44 only for the exact monotone
full-window hierarchy.  It does not apply to optional weight truncations,
arbitrary implemented diagnostic levels, solver-unknown statuses, or
nonmonotone residual lists.

Dependencies: CA-43, CA-44, CA-62, CONVENTIONS.md (n), CONVENTIONS.md (p).

% Source: literature/md/2010.11121/2010.11121.md:126--180 -- inductive-limit
%   states, projective-limit state architecture, and GNS representation.
% Source: literature/md/2010.11121/2010.11121.md:226--243 -- vacuum invariance
%   and warning that Poincare covariance is additional.
% Source: references/qft/FewsterRejzner2019/source/AQFTIMPRS-Oct2019.tex:337--447 --
%   states and representations of unital *-algebras and the GNS representation
%   theorem with its construction steps.
% Source: references/qft/FewsterRejzner2019/source/AQFTIMPRS-Oct2019.tex:318--323 --
%   density matrices induce states; not all algebraic states arise this way.
% GAP: needs local source for Banach-Alaoglu compactness of the state space of
%   a unital C*-algebra.
% GAP: needs local source for product-trace conditional expectations onto
%   finite tensor factors of the spin-chain UHF algebra.

\subsection{Canonical Full-Window Levels}

Let
\[
  I_N=[-N,N],\qquad
  \Alg_N=\bigotimes_{j\in I_N}M_2(\mathbb C),
\]
and let \(\mathcal W_N\) be the set of all Pauli words supported in \(I_N\).
For this shard, a level means the full-window problem:
\[
  y_\emptyset=1,\qquad y_s=y_{\tau^k s}
  \tag{63.1}
\]
whenever both words fit in the level window,
\[
  M_N(y)_{u,v}=\omega(P_u^*P_v)\succeq0,
  \qquad u,v\in\mathcal W_N,
  \tag{63.2}
\]
and every fixed residual relation
\[
  \omega(X^*RY)=0
  \tag{63.3}
\]
whose product support fits in \(I_N\).  The residual family is fixed in
advance, or is enlarged monotonically.  This monotonicity is part of the
theorem; without it the compactness argument below has no nesting premise.

\subsection{Finite Moments Are Finite States}

Given a feasible \(y\) at level \(N\), define a linear functional on
\(\Alg_N\) by
\[
  \omega_N(P_s)=y_s
\]
and extend linearly in the Pauli-word basis.  If
\(X=\sum_{u\in\mathcal W_N}c_uP_u\), then
\[
  \omega_N(X^*X)
  =
  \sum_{u,v}\overline{c_u}c_v\,\omega_N(P_u^*P_v)
  =
  c^*M_N(y)c
  \ge0.
  \tag{63.4}
\]
Since every positive element in a finite matrix algebra has a square root,
(63.4) makes \(\omega_N\) positive.  The normalization \(y_\emptyset=1\)
makes it unital, hence a state on \(\Alg_N\).

Equivalently, with \(d_N=2^{|I_N|}\),
\[
  \rho_N=d_N^{-1}\sum_{s\in\mathcal W_N}y_sP_s
  \tag{63.5}
\]
satisfies \(\omega_N(A)=\Tr(\rho_NA)\), the density-matrix presentation of the
finite state.  Thus \(M_N(y)\succeq0\) if and only if \(\rho_N\succeq0\).

\subsection{Compactness}

Let \(\Alg\) be the quasi-local spin C*-algebra, the norm closure of the
local algebras.  The normalized product trace outside \(I_N\) gives a unital
completely positive conditional expectation
\[
  E_N:\Alg\to\Alg_N.
\]
Therefore \(\omega_N\circ E_N\) is a global state and agrees with
\(\omega_N\) on \(\Alg_N\).  This proves nonemptiness of the level set
below whenever the level moment problem is feasible.

Let \(K_N\subset S(\Alg)\) be the set of global states whose restriction
to \(\Alg_N\) satisfies the affine level-\(N\) translation and residual
constraints.  Positivity is automatic for elements of the state space.  Each
constraint is a finite linear equation in evaluations
\(\varphi\mapsto\varphi(A)\), hence \(K_N\) is weak-* closed.  Since
\(S(\Alg)\) is weak-* compact, each \(K_N\) is compact.

The full-window levels are nested.  Every level-\(N\) translation equality and
every level-\(N\) residual test also appears at level \(N+1\).  Hence
\[
  K_{N+1}\subseteq K_N.
  \tag{63.6}
\]
If every full-window level is feasible, the \(K_N\) are nested nonempty compact
sets, so
\[
  \bigcap_{N\ge0}K_N\neq\emptyset.
  \tag{63.7}
\]
Choose \(\omega\) in this intersection.

For any Pauli word \(s\) and any integer shift \(k\), choose \(M\) so large
that the intermediate translates from \(s\) to \(\tau^k s\) fit in \(I_M\).
Since \(\omega\in K_M\), (63.1) gives
\[
  \omega(P_s)=\omega(P_{\tau^k s}).
\]
By linearity, density of the local Pauli algebra, and continuity of the shift
automorphism, \(\omega\circ\tau^k=\omega\) on all of \(\Alg\).

\subsection{Residual Relations in the GNS Representation}

Let \((\mathcal H_\omega,\pi_\omega,\Omega_\omega)\) be the GNS
representation.  Fix one residual density \(R\).  For local Pauli words
\(X,Y\), choose \(N\) so that \(X^*RY\) fits the level window.  Because
\(\omega\in K_N\),
\[
  \langle \pi_\omega(X)\Omega_\omega,
          \pi_\omega(R)\pi_\omega(Y)\Omega_\omega\rangle
  =
  \omega(X^*RY)
  =
  0.
  \tag{63.8}
\]
The vectors \(\pi_\omega(A)\Omega_\omega\), with \(A\) local, are dense in the
GNS Hilbert space.  Since \(\pi_\omega(R)\) is bounded, (63.8) on this dense
subspace implies
\[
  \pi_\omega(R)=0.
  \tag{63.9}
\]
Translation-invariant states implement the lattice shift by a GNS unitary, so
the translated residual densities vanish as well.

\subsection{Exact Alternative}

Conversely, if a translation-invariant state \(\omega\) satisfies all residual
relations, then \(y_s=\omega(P_s)\) is feasible at every full-window level.
Indeed, \(y_\emptyset=1\), translation invariance gives (63.1), positivity of
\(\omega\) gives (63.2), and the residual hypothesis gives (63.3).

Thus, for fixed exact residual data, there is no third mathematical
possibility:
\[
  \text{some full-window level is infeasible}
  \quad\text{or}\quad
  \text{a global reference state exists.}
  \tag{63.10}
\]
This is an exact compactness statement, not a finite-time algorithm.

\subsection{Non-Implications}

The state in (63.7) need not be unique.  It need not be a ground state, a
positive-energy state, a clustering state, or a scaling-limit state.  The
argument does not construct a continuum limit and does not produce a
projective-unitary Poincare or Virasoro representation.  It only proves that
the imposed finite residual relations have a genuine C*-algebraic state model.

\subsection{Boundary}

Finite feasibility at checked levels still proves no existence theorem.  The
compactness result requires exact feasibility at every monotone full-window
level.  Numerically, exclusion remains semi-decidable by finding a certified
infeasible finite level; existence is not certified by any finite prefix of
feasible solver runs.  The finite witnesses for the proof ingredients
(density-matrix reconstruction, extension moment preservation, level nesting,
and a positivity mutation) are checked in the Julia testset
``qubit full-window moment compactness witnesses''.
